% Draft: selection of positive TSS classes (P+S)
% Target: Chapter 4 (Materials and Methods), dataset reconstruction;
%         short justification in Chapter 7 (Experimental Protocol).
% Status: draft v1, 2026-11-02. References checked against Crossref DOIs.
% Project sources: specs/graphpro_dataset_reconstruction.md (TSS classes),
%   docs/provenance/research_log.md (R-035), findings F-115 to F-117.

\section{Selection of Positive TSS Classes}\label{draft:tss-class-selection}

\draftmeta{Chapter~4 (dataset reconstruction); summary in Chapter~7}
{Draft v1 --- references verified; positive-class rule awaiting owner approval}
{\texttt{specs/graphpro\_dataset\_reconstruction.md}; log R-035; F-115--F-117}

\subsection{TSS maps are not promoter lists}

\citet{sharma2010primary} introduced differential RNA-seq (dRNA-seq) to map the
primary transcriptome of \textit{Helicobacter pylori} and classified every
transcription start site (TSS) by its position relative to annotated genes:
\begin{itemize}
  \item \textbf{primary (P):} the main TSS upstream of a gene, on the same
    strand;
  \item \textbf{secondary (S):} additional TSSs upstream of the same gene,
    with weaker support;
  \item \textbf{internal (I):} inside a coding sequence, on the sense strand;
  \item \textbf{antisense (A):} on the opposite strand, inside or near a gene;
  \item \textbf{orphan (O):} not associated with any annotated gene.
\end{itemize}
The classes are not mutually exclusive. Later dRNA-seq studies adopted the
scheme with study-specific distance thresholds; \citet{dugar2013campylobacter},
for example, assign a TSS to P or S when it lies at most 300~bp upstream of a
gene, and to A when it lies inside a gene or within 100~bp on the opposite
strand. Maps built this way exist for \textit{Escherichia coli}
\citep{thomason2015ecoli}, \textit{Campylobacter jejuni}
\citep{dugar2013campylobacter}, \textit{Chlamydia pneumoniae}
\citep{albrecht2011chlamydia}, and \textit{Leptospira interrogans}
\citep{zhukova2017leptospira}.

A dRNA-seq map is therefore an inventory of heterogeneous transcription
initiation events, not a list of promoters. In \textit{E.~coli}, 5,574 of
14,868 TSS candidates were internal to annotated genes and 5,495 corresponded
to potential antisense transcripts \citep{thomason2015ecoli}. The set of
classes accepted as positive is part of the definition of the prediction task:
two datasets built from the same organism and experiment but with different
class sets measure different things. The protocol therefore declares the
positive class set for each organism before any dataset is built or any model
is trained, and never tunes it against model performance.

\subsection{Rationale for P+S}

\paragraph{Construct validity.}
In the canonical molecular definition, a promoter is the DNA region where RNA
polymerase, guided by a $\sigma$ factor, initiates and regulates the
transcription of a gene \citep{browning2016regulation}. P and S are exactly the
classes defined by this relationship: they lie upstream of a gene, on its
strand. Restricting positives to P+S makes the positive label denote one
biological construct rather than the mixture of initiation events in a
complete map.

\paragraph{Homogeneity of the signal and avoidance of compositional shortcuts.}
Internal TSSs, and many antisense TSSs, lie inside protein-coding sequences.
Their windows inherit coding-sequence statistics, such as codon periodicity and
position-specific base composition, which identify coding regions on their own
\citep{fickett1982coding}. Mixed into the positive class, they would let a
classifier separate the classes by learning ``coding versus non-coding''
instead of ``promoter versus non-promoter''. This is a shortcut: a decision
rule that succeeds on the benchmark but does not transfer to the intended task
\citep{geirhos2020shortcut}. It would also contradict any negative condition
derived from coding sequences, where the same kind of sequence would carry the
opposite label. P+S keeps the positives in the upstream region where the
canonical promoter elements reside \citep{browning2016regulation}.

\paragraph{Strength and stability of the evidence.}
Antisense and orphan TSSs include a large share of pervasive, low-level
transcription whose function is debated \citep{wade2014pervasive}. Across 20
bacterial species, the number of antisense RNAs increases with genomic AT
content, and most antisense transcripts are consistent with initiation at
spurious promoters, that is, transcriptional noise
\citep{llorensrico2016antisense}. Gene-associated TSSs are also the most
reproducible across strains \citep{dugar2013campylobacter} and growth
conditions \citep{zhukova2017leptospira}. Anchoring the positive class in
gene-associated TSSs reduces label noise without requiring any judgement on
whether the other classes are biologically functional.

\paragraph{Why S is included and the class is not restricted to P.}
The distinction between P and S ranks the TSSs of one gene by the signal
measured under the conditions of the experiment; it is not a biological
category. Genes are frequently transcribed from several promoters, used
differently across conditions and $\sigma$ factors
\citep{cho2009transcription,mendozavargas2009genome}, so a secondary TSS in one
condition may be the dominant TSS in another. A P-only definition would encode
an artefact of the experimental condition and concentrate the positive class on
the strongest promoters. Evaluating on such a narrow spectrum of cases tends to
overstate performance relative to the full population, as described for
diagnostic tests \citep{ransohoff1978spectrum}. Including S preserves the
natural diversity of gene-associated promoters and increases the sample size
without mixing constructs.

\paragraph{Comparability across organisms and sources.}
Because P and S are defined relative to the annotated gene, they are the classes
most consistently present across dRNA-seq studies that follow the scheme of
\citet{sharma2010primary}
\citep{albrecht2011chlamydia,dugar2013campylobacter,thomason2015ecoli,zhukova2017leptospira}.
Applying one rule to each primary table yields datasets whose labels have the
same meaning across organisms. This is a precondition for cross-dataset
inference: comparing effects across datasets is meaningful only if the positive
label denotes the same construct in each. The point also applies to released
benchmarks; the Promotech training inventory
\citep{chevezguardado2021promotech} applies no common class rule across
organisms (finding F-115).

\paragraph{Coherence with negative construction.}
TSSs outside the selected classes are not used as negatives. Genomic controls
are excluded from the neighbourhood of every TSS in the complete primary table,
for all classes and on both strands, so no control contains a known initiation
site, and the I, A, and O events are removed from both classes instead of
contaminating the negative class. This is what makes restricting the positive
class defensible. Random genomic controls remain unlabelled rather than
verified negatives, and this residual uncertainty is reported as a
positive--unlabelled setting \citep{elkan2008pu}.

\subsection{Operational rules}

\begin{enumerate}
  \item A TSS is positive if it holds at least one of the classes P or S,
    regardless of any other class it also holds.
  \item Classes are taken from the primary table of the source study, never
    from derived subsets.
  \item The full class set of every positive is retained in the pair ledger,
    so results can be stratified by class.
  \item Alternative class sets (P only, P+S+I, all classes) are run only as
    separate, versioned sensitivity conditions and are never pooled with the
    primary condition.
  \item Class assignment depends on the genome annotation used by the source
    study; that annotation version is recorded as provenance.
\end{enumerate}

\subsection{Limitations}

P+S inherits annotation errors: a promoter of an unannotated gene is
classified as orphan and excluded. The definition also excludes genuine
promoters of antisense and intragenic transcripts, some of which are functional
\citep{wade2014pervasive}; the task is therefore explicitly gene-associated
promoter prediction, not prediction of all transcription initiation. When a
source releases only primary TSSs, P+S reduces to P, and this asymmetry between
sources must be reported. Distance thresholds for P/S assignment differ between
studies \citep[e.g.][]{dugar2013campylobacter}, so the threshold applied by each
source is recorded.
